\cotefile{156-6}
\scanpage{18}
\mspage{17}

Il n'y a pas à y \unsure{revenir}. Je vais plutôt m'attacher à présent aux interprétations en termes d'ombres (grandes ou petites). Il est entendu que dans un premier temps, je m'attache aux \unsure{aspects|aspect} ``algébriques'' ou ``triviaux'' d'un formalisme des figures, et me contente de dégager les bons axiomes, de telle façon que

1\textsuperscript{o}) Ceux-ci \unsure{soient|doivent} valables dans tout ``\unsure{atelier|atlas}'' de la forme $\mathrm{Fig}(\deleted{\unsure{E|L}}L)$, formé des figures \unsure{ensemblistes} dans un ensemble donné $L$ (sans conditions de finitude d'aucune sorte),

et

2\textsuperscript{o}) qu'ils soient applicables aussi, disons, aux figures \uline{simpliciales} affines finies, dans un espace affine sur $\mathbb{R}$ (ou, plus \unsure{gén\textsuperscript{t}.|gént.}, sur un corps ordonné quelconque).

Pour expliciter, en termes de l'ens ordonné $(M,\leq)$, \unsure{l'ens|les ens} $\unsure{\mathcal{F}|\mathscr{F}}$ des figures, en tant qu'ens de parties de $M$, \deleted{\unsure{il est}} ou du moins des \uline{figures de type fini}
\[
\begin{aligned}
 F\in\mathcal{F}\ \text{de t.f} \quad&\Longleftrightarrow\quad
 \widetilde F=\deleted{\unsure{\text{fibres}|\text{fibre}}}\,\text{réunion finie}\\
 &\hspace{4.9em}\text{d'\unsure{un|ens}}\,M_{\leq x_i}\\
 &\hspace{3.0em}\text{(pour un nb \uline{fini}}\\
 &\hspace{3.5em}\text{de strates de }\unsure{M|x_i|X}\text{)}
\end{aligned}
\tag*{(1.34)}
\]
\scanpage{19}
\mspage{18}

la chose essentielle est la connaissance de
\[
\underset{\unsure{M|m}}{\mathrel{\tikz[baseline=-.55ex,x=1em,y=1em]{\draw[line width=.065em,line cap=round,line join=round] (0,.40)--(.48,0)--(0,-.35) (.51,.40)--(.03,.02)--(.49,-.35);}}}\quad\text{relation binaire sym. réfl. (sur }M\text{)}
\tag*{(1.35)}
\]
\deleted{\unsure{qui indique}} satisfaisant
\[
 \unsure{X|x} \mathrel{\tikz[baseline=-.55ex,x=1em,y=1em]{\draw[line width=.065em,line cap=round,line join=round] (0,.40)--(.48,0)--(0,-.35) (.51,.40)--(.03,.02)--(.49,-.35);}} \unsure{Y|y},\quad X'\leq X,\,Y'\leq Y\quad\Longrightarrow\quad X' \mathrel{\tikz[baseline=-.55ex,x=1em,y=1em]{\draw[line width=.065em,line cap=round,line join=round] (0,.40)--(.48,0)--(0,-.35) (.51,.40)--(.03,.02)--(.49,-.35);}} \unsure{Y'|Y,}
 \tag*{(1.35)}
\]

\deleted{\unsure{elle|alors}}

il y a corr. biuniv. \unsure{entre|en} l'ens \unsure{des|de} telles relations sur $M$, et l'ens des $\unsure{\mathcal{F}|\mathscr{F}}\inserted{\subset\unsure{\mathcal{P}|P}(M)}$ satisfaisant les conditions énumérées précédemment, plus celle\deleted{\unsure{s} d'être} que les $F\in\mathcal{F}$ soient \deleted{\unsure{toutes|toute}} de type fini\unsure{.|e} C'est en ce sens que j'ai pu dire que ``\unsure{moralement|essentiellement}'', le couple $(M,\mathcal{F}\subset\mathcal{P}(M))$ équivalait à \deleted{\illegible{1}} la donnée de $\underset{M}{\leqslant}$ \unsure{seul|et|soit} \unsure{de|des} $\mathrel{\tikz[baseline=-.55ex,x=1em,y=1em]{\draw[line width=.065em,line cap=round,line join=round] (0,.40)--(.48,0)--(0,-.35) (.51,.40)--(.03,.02)--(.49,-.35);}}$ satisfaisant (1.36).

\deleted{Mais il s'agit maintenant}

Pratiquement, dans presque tous les exemples, c'est via $M$\inserted{(\unsure{multistrates|multistrate})}, muni de $\leq$ et d'une relation de disjonction $\mathrel{\tikz[baseline=-.55ex,x=1em,y=1em]{\draw[line width=.065em,line cap=round,line join=round] (0,.40)--(.48,0)--(0,-.35) (.51,.40)--(.03,.02)--(.49,-.35);}}$, qu'on \unsure{définira|définirait} les ateliers\ldots

Il faut à présent en venir à la relation $\ll$. \deleted{\unsure{Comme}} \unsure{Elle} va pour l'instant \unsure{servir} surtout à former des ombres\unsure{,|.} Je commence par les grandes ombres
\scanpage{20}
\mspage{19}

\[
 \mathrm{Omb}(F)=\{X\in M\mid X\ll F\}\,\deleted{\unsure{\subset|=}}
 \tag*{$(*)$}
\]

et la multiombre
\[
 \unsure{\mathrm{Multom}|\mathrm{Multiom}|\mathrm{Multiomb}}(F)=\{\mathrm{Omb}(X)\mid X\in\widetilde F\}.
 \tag*{$(**)$}
\]
% Right margin, with a leader to the tilde on F.
\margin{$\widetilde F=F$, si on identifie $F$ à une partie de $M$ via $\unsure{\mathcal{F}|\mathscr{F}|\mathbb{F}}\unsure{\to|\mapsto}\unsure{\mathcal{P}|P}(M)$\\ $F\mapsto\widetilde F$}

En apparence $(*)$ utilise la connaissance de $\ll$ dans $\mathcal{F}$ tout entier, mais \deleted{\illegible{1}}\inserted{il} \deleted{\unsure{on dira qu'elle}} \inserted{résultera des axiomes} qu'elle ne dépend que de $\ll$. En tous cas, il est clair que la multiombre ne dépend que de $\underset{\unsure{M|m}}{\ll}$ (au lieu de $\underset{\mathcal{F}}{\ll}$).

% slanted
\margin{\uline{NB} Je vais changer plus bas le point de vue \unsure{pris|pour} $\ll$ et les axiomes C4, C5.}

Je veux que cette multiombre soit une \uline{figure ensembliste dans} $M$, \uline{ayant pour support l'ombre}, \inserted{et indexée fidèlement par $\widetilde F$\,$\mathord{\tikz[baseline=-.2ex,x=1em,y=1em]{\draw[line width=.055em] (0,0) circle (.13) (-.04,-.07)--(.13,.54);}}$} Cela nous amène à
\margin{$\mathord{\tikz[baseline=-.2ex,x=1em,y=1em]{\draw[line width=.055em] (0,0) circle (.13) (-.04,-.07)--(.13,.54);}}$ avec la relation d'ordre $\leq$ dans $\widetilde F$ correspondant à l'inclusion dans $\mathrm{Multiomb}(F)$, i.e. $\downarrow$}

\inserted{\uline{C4} Si $X,Y\in M$ tels que $X \mathrel{\tikz[baseline=-.55ex,x=1em,y=1em]{\draw[line width=.065em,line cap=round,line join=round] (0,.40)--(.48,0)--(0,-.35) (.51,.40)--(.03,.02)--(.49,-.35);}} Y$ (i.e. $X,Y$ font partie d'une \unsure{=|même} figure $F$), alors $X\ll Y\Rightarrow X\leq Y$}

\uline{C\deleted{\unsure{4}}5}. Si $X\in M$, $F\in\mathcal{F}$, \inserted{tels que $X\ll F$, et} si $F^X$ est l'ens
\[
 \text{\deleted{des} des }Y\in M\quad\text{tels que }\deleted{\unsure{X\ll Y|X\leq Y}}\,X\ll Y\leq F
 \tag*{(1.37)}
\]
(\unsure{l'ens|l'une} des strates de $F$ \unsure{raffinées|raffinée} par $X$), \deleted{\unsure{on}} cet ens $F^X$ \unsure{.|,} admet un plus petit élément \inserted{pour $\leq$} (i.e. il existe une strate de $F$ raffinée par $X$, et parmi \deleted{\unsure{toutes}} les strates de $F$ raffinées par $X$, il en est une plus petite)

% slanted
\margin{\uline{Cor} \inserted{Si $X\in M$, $F\in\mathcal{F}$} Pour que $X\ll F$, il \unsure{f et S|f et s} qu'il existe $Y\deleted{\in\widetilde F}\in M$ avec $X\ll\underbrace{Y\leq F}_{Y\in\widetilde F}$. Ainsi $X\ll F$ \unsure{de|dans} $M\times\mathcal{F}$ connu via $X\ll Y$ \unsure{de|dans} $M\times M$}

\deleted{\uline{C6} Soit $\unsure{F\in\mathcal{F},\,X\in\widetilde F}$ \unsure{soit|puis} $X\in M$, $\unsure{X\ll F|X\in\mathcal{F}}$.}

\deleted{Alors $\exists Y\in\unsure{M|\mathcal{F}M}$ avec \unsure{*} a) $Y\ll X$ i.e. $X\in F^Y$ \quad b) $Y\ll X'\leq X$ et $X'\in M\Longrightarrow X'=X$}

\deleted{b) \unsure{que|pour} $X$ est la plus \uline{petite} strate de \unsure{$F$} (cf C4)}

\margin{\uline{NB} Le fait que dans $\mathrm{Multom}(X)$, les $\mathrm{Omb}(X)^{\unsure{\circ|o}}\ne\varnothing$, est trivial\,; car $X\in\mathrm{Omb}(X)^\circ$}
\scanpage{21}
\mspage{20}

Ainsi on trouve donc une application
\[
 \unsure{\mathcal{F}|\mathscr{F}}\longrightarrow\mathrm{Fig}(M),\quad F\longmapsto\unsure{\mathrm{Multomb}|\mathrm{Multiomb}}(F)
 \tag*{$(*)$}
\]
Elle est injective, car \inserted{$X$ est connu quand $\mathrm{Omb}(X)$ est connu (c'est le plus grand élément de $\mathrm{Omb}(X)$ pour $\ll$), donc $\widetilde F$ est connu quand on \unsure{connaît|connait} $\unsure{\mathrm{Multom}|\mathrm{Multiomb}}(F)$, donc $F$ aussi.}

Il est immédiat \unsure{que}
\[
 \deleted{F}G\leqslant\deleted{G}F\quad
 \unsure{\Longleftrightarrow|\Longrightarrow}\quad
 \underbrace{
 \underbrace{\mathrm{Multomb}(\deleted{F}G)}_{
 \substack{\text{les strates sont}\\\text{les }\mathrm{Omb}(\deleted{X}Y)\\\deleted{X}Y\in\widetilde{\deleted{F}G}}}
 \leqslant
 \underbrace{\mathrm{Multomb}(\deleted{G}F)}_{
 \substack{\text{les strates}\\\text{sont les}\\\mathrm{Omb}(X)\\X\in\widetilde F}}}_{
 \substack{\text{or }\widetilde{\unsure{G|C}}\subset\widetilde F\text{ est}\\
 \text{une partie}\\\text{fermée de }\widetilde F\\\text{si }G\leq F}}
 \tag*{(1.38)}
\]

\uline{19 \unsure{Juin|juin}}\qquad Finalement, j'ai changé de point de vue pour le sens géométrique de la relation $\deleted{F}G\ll\deleted{G}F$, que je \unsure{concevais} comme ``$G$ est une \unsure{sous figure|sousfigure} d'une subdivision de $F$''. Déjà dans \deleted{\unsure{les}} le cas simplicial (i.e la structure \unsure{combinatoire} des \deleted{\unsure{tout|tous}} multistrates étant \deleted{celle} une structure de simplexe : \unsure{l'ens.|ens.} ordonné des sous-ens d'un ens. fini \ldots), cela n'a pas l'air naturel. Par exemple, si on a une théorie modérée, et deux simplexes modérés $X,Y$ dans un $\mathbb{R}^{\unsure{n}}$ disons, on dira que $\deleted{X}Y\ll X$ si c'est le cas pour les figures ensemblistes correspondantes -- ce qui n'implique sans doute pas l'existence d'une \unsure{subd.|subdiv.} de $X$ telle que $Y$ en soit une multistrate \unsure{et|de} la
\scanpage{22}
\mspage{21}

question s'il en existe une semblerait \unsure{accessoire|nécessaire}. Aussi, finalement, j'ai envie de poser

\noindent\begin{tabular}{@{}||@{\hspace{.55em}}p{.93\linewidth}@{}}
\uline{A\deleted{C}\unsure{t|l|x} 4}\quad Si $\deleted{G}F,\deleted{F}G\in\unsure{\mathcal{F}|\mathscr{F}}$, alors $\deleted{F}G\ll\deleted{G}F$ sss pour tout $X\in\widetilde{\deleted{F}G}$, $\deleted{\unsure{\exists|s}}$ $\exists Y\in\widetilde F$ avec $X\ll Y$
\end{tabular}

Ainsi, la relation $\ll$ sur $\mathcal{F}$ se déduit de la relation $\ll$ sur $M$. Donc, si on \unsure{suppose|s'impose} les figures de type fini, $(\mathcal{F},\leq,\ll)$ (i.e. tout l'atelier) se réconstitue entièrement à partir de $(M,\leq,\ll,\mathrel{\tikz[baseline=-.55ex,x=1em,y=1em]{\draw[line width=.065em,line cap=round,line join=round] (0,.40)--(.48,0)--(0,-.35) (.51,.40)--(.03,.02)--(.49,-.35);}})$.
% slanted
\margin{\uline{Cor} $\forall F\in\mathcal{F},\ \mathrm{Omb}(F)=\bigl|\unsure{\mathrm{Multomb}|\mathrm{Multiomb}}\,F\bigr|$}

\noindent\begin{tabular}{@{}||@{\hspace{.55em}}p{.93\linewidth}@{}}
\uline{\deleted{C}At 5}\quad Soient $X\in M$, $F\in\mathcal{F}$, tels que $X\ll F$. Alors \deleted{dans} l'ensemble $F^X=\{Y\in\widetilde F\mid\underbrace{X\ll Y}_{\text{i.e. }X\in\mathrm{Omb}(Y)}\}$ (lequel est $\ne\varnothing$ par C4) admet un plus petit élément\unsure{,|.|*} (pour $\leq$)
\end{tabular}

En d'autres termes, l'ensemble de parties $\mathrm{Multom}\deleted{\unsure{b}}b(F)=\{\widetilde Y\mid Y\in\widetilde F\ \text{i.e.}\ Y\triangleleft F\}$ est une \uline{figure} dans $M$, \deleted{[\uline{NB} \unsure{étant entendu} que les $\widetilde Y^{\unsure{\circ|o}}=\widetilde Y\setminus\bigcup_{\unsure{X<|X\ll}}$} ou du moins, est \uuline{pré}\uline{admissible}. C'est en tous cas impliqué par \uline{C5}, et lui est équivalent si on admet que la relation (pour $Y,Y'\in\widetilde F$) $Y\leq Y'$, qui implique $Y\ll Y'$ \unsure{d'où|donc} $\deleted{\unsure{G|O}}\mathrm{Omb}(Y)\subset\mathrm{Omb}(Y')$ lui est en fait équivalente. Cela équivaut à la condition

\noindent\begin{tabular}{@{}||@{\hspace{.55em}}p{.93\linewidth}@{}}
\uline{\deleted{C}At 6}\quad Si $X,Y\in M$ et $X \mathrel{\tikz[baseline=-.55ex,x=1em,y=1em]{\draw[line width=.065em,line cap=round,line join=round] (0,.40)--(.48,0)--(0,-.35) (.51,.40)--(.03,.02)--(.49,-.35);}} Y$, alors $X\ll Y\Rightarrow X\leq Y$ (donc $X\ll Y\ \unsure{\Longleftrightarrow|\longleftrightarrow}\ X\leq Y$).

En d'autres termes,
\end{tabular}
\scanpage{23}
\mspage{22}

\uline{Corollaire 1} Soit $F\in\unsure{\mathcal{F}|\mathscr{F}}$. Alors dans $\widetilde F\subset M$, les relations d'ordre $\leq$ et $\ll$ coïncident.

\uline{Corollaire 2} Soient $F$ et $\unsure{G|C}$ deux figures compatibles $(F \mathrel{\tikz[baseline=-.55ex,x=1em,y=1em]{\draw[line width=.065em,line cap=round,line join=round] (0,.40)--(.48,0)--(0,-.35) (.51,.40)--(.03,.02)--(.49,-.35);}} G)$. Alors $F\ll G$ sss $F\leq G$.

Il suffit de prouver $\Longrightarrow$, et par C4, on est ramené à \deleted{\unsure{prouver que}}\inserted{le voir} pour $X$ ($\in\widetilde F$) et $Y$ ($\in\widetilde G$) \unsure{dans|sous} $M$, ce qui n'est autre que C6.

Ainsi, pour tout $F\in\mathcal{F}$, l'indexation de $\mathrm{Multomb}(F)$ par $\widetilde F$ est fidèle, et respecte les relations d'ordre naturelles \inserted{\unsure{habituelle|habituelles} $\leq$}.

Ceci \deleted{\unsure{a|e}} dit, pour voir que $\mathrm{Multomb}(F)$ est bien une figure, il reste à voir que
\[
 \text{pour }Y\in\widetilde F,\qquad \widetilde Y^{\unsure{\circ|o}}\overset{\mathrm{def}}{=}
 \widetilde Y\setminus
 \bigcup_{\substack{Y'\in\widetilde F\\Y'<Y\;(\text{ou encore, }Y'\ll Y,\,Y'\ne Y)}}\widetilde Y'\quad\text{est}
\]
non vide. Mais il n'y a qu'à \deleted{\unsure{prouver|peut-être|remarquer}}\inserted{noter} que $Y\in\widetilde Y^\circ$ !
% slanted
\margin{NB. \unsure{pourrait} remonter avant C6, car il n'y avait pas besoin de savoir que $Y'\ll Y$, $Y'\ne Y$ équivaut à $Y'<Y$ pour $Y,Y'\in\widetilde F$ !}

Par la relation d'ordre $\ll$ sur $M$, on trouve une application
\[
 \begin{aligned}
 M&\lhook\joinrel\longrightarrow\unsure{\mathcal{P}|P}(M)\\
 X&\longmapsto\mathrm{Omb}(X)=M_{\ll X}=\{Y\in M\mid Y\ll X\}
 \end{aligned}
 \tag*{(\deleted{\unsure{*|1.38}}1.39)}
\]
qui est \uline{injection} ($\ll$, relation \uline{d'ordre}, pas seulement préordre) \inserted{et la relation d'ordre de $M$ \unsure{induit|induite} par celle de $\mathcal{P}$}. Les figures $\unsure{\mathrm{Multom}|\mathrm{Multomb}}(F)$ sont les images, par cette application, des $\widetilde F$ ($F\in\mathcal{F}$). \deleted{Si $F$ et $G\in\mathcal{F}$,}
\scanpage{24}
\mspage{23}

\uline{Proposition} Soient $F,\unsure{G|C}\in\unsure{\mathcal F|\mathscr F}$. Conditions équivalentes
\[
\begin{aligned}
&\text{(i)}\quad F\leq G\\
&\text{(i bis)}\quad \unsure{\widetilde F|\widetilde{\widetilde F}}\subset\widetilde G\\
&\text{(ii)}\quad \unsure{\mathrm{Multomb}|\mathrm{Multiomb}}(F)\leq\mathrm{Multomb}(G)
 \qquad\begin{gathered}\text{i.e. }\mathrm{Multomb}(F)\\\text{est une sous-figure}\\\text{de }\mathrm{Multomb}(G)\end{gathered}\\
&\text{(ii bis)}\quad \mathrm{Multomb}(F)\subset\mathrm{Multomb}\unsure{G)|(G)}
\end{aligned}
\]

Dém. $F\leq G$ \deleted{trivial, et} $\unsure{\Longleftrightarrow|\Longleftarrow}\widetilde F\subset\widetilde G$ \deleted{trivial, et}\inserted{déjà connu} \deleted{\unsure{l'inverse}}\inserted{et mis} pour mémoire ($\Longrightarrow$ trivial, $\Longleftarrow$ car $\deleted{\widetilde F}F=\mathop{\mathrm{Sup}}\limits_{X\in\widetilde F}X$, $G=\mathop{\mathrm{Sup}}\limits_{Y\in\widetilde G}Y$). L'équivalence (i bis) $\Longleftrightarrow$ (ii bis) triviale par l'injectivité de (\deleted{\unsure{*|1.38}}1.39). Or (ii bis) est impliqué par (ii) (tautologique) prouvons l'inverse, donc que (ii bis) implique que $\unsure{\mathrm{Mulomb}|\mathrm{Multomb}}(F)$ est fermé dans $\mathrm{Mulom}\,G$ pour l'inclusion. Mais on aura $\widetilde F\subset\widetilde G$ \deleted{et il est}\inserted{donc} \deleted{\unsure{et on|et en}} $F\leq G$, et il est \deleted{\unsure{connu}}\inserted{connu et trivial} que $\widetilde F$ est \inserted{alors} fermé dans $\widetilde G$.

\uline{Corollaire 1} \deleted{Conditions équivalentes pour}\inserted{Soient} $F,G\in M$ \inserted{tels que $F \mathrel{\tikz[baseline=-.55ex,x=1em,y=1em]{\draw[line width=.065em,line cap=round,line join=round] (0,.40)--(.48,0)--(0,-.35) (.51,.40)--(.03,.02)--(.49,-.35);}} G$, alors} \deleted{(i) $F \mathrel{\tikz[baseline=-.55ex,x=1em,y=1em]{\draw[line width=.065em,line cap=round,line join=round] (0,.40)--(.48,0)--(0,-.35) (.51,.40)--(.03,.02)--(.49,-.35);}} G$}
\[
 \deleted{\text{(ii)}}\quad \mathrm{Mulomb}(F)\ \mathrel{\tikz[baseline=-.55ex,x=1em,y=1em]{\draw[line width=.065em,line cap=round,line join=round] (0,.40)--(.48,0)--(0,-.35) (.51,.40)--(.03,.02)--(.49,-.35);}}\ \mathrm{Mulomb}(G),
 \tag*{(*)}
\]
\deleted{Quand elles sont remplies,}\inserted{et} on a
\[
 \left\{\begin{aligned}
  \mathrm{Mulomb}(F\vee G)&=\mathrm{Mulom}(F)\cup\mathrm{Mulomb}(G)\\
  \mathrm{Mulomb}(F\wedge G)&=\mathrm{Mulomb}(F)\cap\mathrm{Mulomb}(G).
 \end{aligned}\right.
 \tag*{(**)}
\]

\deleted{La condition (ii) signifie que $\widetilde F$ et $\widetilde G$ sont fermés dans $\widetilde F\cup\widetilde G\subset M$\unsure{,|\ensuremath{{}'}} pour la relation $\ll$ (induite par $\subset$ dans $\unsure{\mathcal P|P}(M)$).}
\scanpage{25}
\mspage{24}

En effet, $\widetilde{F\vee\unsure{G|C}}=\widetilde F\cup\widetilde G$, et $\widetilde F$ et $\widetilde G$ y sont fermés, \deleted{dans $M$} ce qui prouve que
\[
 \unsure{\mathrm{Multomb}|\mathrm{Multiomb}}(F\unsure{\cup|\vee}G)=\mathrm{Multomb}(F)\cup\mathrm{Mulomb}(G),
\]
et que $\mathrm{Mulomb}(F)$ et $\mathrm{Mulomb}(G)$ y sont fermés, donc \unsure{(*)} et la première relation (**). La deuxième relation \inserted{(**)} provient de $=$ de $\widetilde{F\cap G}=\widetilde F\cap\widetilde G$.

\uline{NB} Il n'est pas vrai en général que $\mathrm{Mulomb}(F) \mathrel{\tikz[baseline=-.55ex,x=1em,y=1em]{\draw[line width=.065em,line cap=round,line join=round] (0,.40)--(.48,0)--(0,-.35) (.51,.40)--(.03,.02)--(.49,-.35);}} \mathrm{Mulomb}(G)$ implique $F \mathrel{\tikz[baseline=-.55ex,x=1em,y=1em]{\draw[line width=.065em,line cap=round,line join=round] (0,.40)--(.48,0)--(0,-.35) (.51,.40)--(.03,.02)--(.49,-.35);}} G$, même si $F$ et $G$ sont réduits aux lieux distincts $x$ et $y$ --- auquel cas $\mathrm{Multomb}(X)=\{\{\unsure{\lambda|x}\}\}$, $\mathrm{Mulomb}(Y)=\{\{y\}\}$ sont compatibles et $=$ disjointes, mais $x$ et $y$ ne le sont pas forcément. Mais il \unsure{sera|ser.} vrai \unsure{par contre qu'on a équivalence|pour d'autres quasi-équivalences|pour d'autre quasi-équivalence}, dans le cas où $\unsure{\mathcal F|\mathscr F}$ est \unsure{un|une} \unsure{atelier} \unsure{ensembliste|ens.embliste}.
% slanted
\margin{\unsure{vérifier|vérifies} plus tard}

\uline{Corollaire 2} Si $F$ et $G$ sont disjoints, alors $\unsure{\mathrm{Mul.mb}|\mathrm{Mulomb}}(F)$ et $\mathrm{Mulomb}(G)$ aussi.

Cela provient formellement du corollaire 1 et de la formule
\[
 \mathrm{Mulomb}(\varnothing_{\mathcal F})=\varnothing\quad\text{(figure vide dans $M$)}
 \tag*{(1.40)}
\]

\uline{Proposition} \unsure{Soient|Soit} $\unsure{F|\mathbb F},G\in\mathcal F$. Alors $F\ll G$
\[
 \Longleftrightarrow\quad\mathrm{Mulomb}(F)\ll\mathrm{Mulomb}(G).
\]
